\section{核心科学问题}\label{sec:problems}

本研究的学术价值根植于若干具有根本意义的开放科学问题。围绕PrimeCoDec框架所要解决的三大结构性困境，我们将其提炼为以下五个相互关联的核心科学问题：

\vspace{6pt}
\begin{tcolorbox}[problemstyle]
\textbf{科学问题 P1：通用信道统一表示}

如何为来自截然不同物理机制（电磁传播、光子噪声、量子退相干、固态器件噪声）的异质噪声信道，建立一种\textbf{数学完备且计算高效}的统一表示框架，使得基于此框架设计的AI模型能够在不同系统间无缝迁移，而无需针对每类系统重新训练或重新设计？

\textit{挑战所在}：不同物理系统的噪声统计结构差异极大——从高斯分布（AWGN）、混合高斯（Flash）到Pauli噪声（量子）；噪声的时空相关性结构从无关（独立信道）到强相关（突发错误信道）跨越多个数量级；概率密度函数可能不存在解析形式（仅有测量数据）。
\end{tcolorbox}

\vspace{6pt}
\begin{tcolorbox}[problemstyle]
\textbf{科学问题 P2：纠错码性能的高精度神经算子预测}

是否存在一种神经网络架构，能够在\textbf{无需蒙特卡洛仿真}的前提下，对任意给定的码结构$\calC$与信道三元组$\calT$，以\textbf{高精度}（BER曲线误差$< 0.1$~dB）、\textbf{超快速度}（推理时间$< 10$~ms）预测完整的BER/FER性能曲线，并对不同码长、码率、校验矩阵与信道参数的组合具有强泛化能力？

\textit{挑战所在}：该映射本质上是从无限维函数空间到无限维函数空间的算子，传统神经网络仅能映射有限维向量，直接应用存在根本性局限。BER曲线在瀑布区（waterfall region）呈现高度非线性的急降特性，对预测模型的局部精度提出了极高要求；BER值跨越$10^{-1}$到$10^{-10}$多个数量级，需要特殊的损失函数设计。
\end{tcolorbox}

\vspace{6pt}
\begin{tcolorbox}[problemstyle]
\textbf{科学问题 P3：算法-硬件联合性能代理建模}

如何构建一个轻量级代理模型，能够同时以高精度预测ECC\textbf{算法层指标}（BER/FER性能）与\textbf{硬件层指标}（芯片面积$A$、功耗$P$、关键路径时延$T_d$、吞吐量$\Gamma$），从而支持在统一目标函数框架下的算法-硬件协同优化？

\textit{挑战所在}：算法性能与硬件指标的影响因素部分重叠、部分正交。硬件指标高度依赖工艺节点与综合工具设置，如何在代理模型中有效编码工艺参数是一个未解问题。
\end{tcolorbox}

\vspace{6pt}
\begin{tcolorbox}[problemstyle]
\textbf{科学问题 P4：大语言模型驱动的需求形式化}

如何设计一个可靠的LLM推理管道，将工程师提出的\textbf{自然语言模糊需求}（如"需要一个用于5G URLLC的、在城市非视距（NLOS）场景下可靠性达$10^{-9}$的纠错方案"）或标准规范文档中的\textbf{上下文约束}（3GPP TS 38.211等），自动转化为优化引擎可直接处理的\textbf{形式化约束集合}$\calC$？该过程的\textbf{完备性}与\textbf{一致性}如何得到保证？

\textit{挑战所在}：工程需求往往包含隐含假设、领域专有术语与相互竞争的多目标约束。LLM在处理此类结构化工程推理时容易出现幻觉（hallucination）或约束遗漏，需要配套的验证机制。
\end{tcolorbox}

\vspace{6pt}
\begin{tcolorbox}[problemstyle]
\textbf{科学问题 P5：行业标准自动化生成}

如何将PrimeCoDec框架输出的最优码型参数、译码算法与硬件实现结果，\textbf{自动逆向映射}为符合3GPP、IEEE 802.11等工业格式的行业标准草案文档？该自动化标准生成过程的\textbf{合规性}、\textbf{完备性}与\textbf{可读性}如何保证？

\textit{挑战所在}：行业标准文档具有严格的格式规范、术语体系与条款逻辑；将数学优化结果翻译为工程规范语言需要深度领域知识；不同应用场景（无线、光通信、存储）对应不同的标准化组织与文档体系，如何实现统一的标准生成接口是一个未解问题。
\end{tcolorbox}

\vspace{6pt}
上述五个科学问题相互关联、层层递进，共同构成PrimeCoDec的理论核心。P1是整个框架的数学基础，使跨域泛化成为可能；P2是突破仿真瓶颈、消解专家依赖的核心技术；P3将性能预测能力扩展至硬件维度，终结算法-硬件割裂；P4是实现端到端自动化闭环的用户交互接口；P5则将框架的价值延伸至行业标准层，使优化结果直接转化为可落地的工业规范。本提案将系统性地攻克这五个问题，建立完整的理论与工程框架。